\documentclass[11pt,a4paper]{amsart}

\usepackage[margin=1in]{geometry}
\usepackage{amsmath,amssymb,amsthm,mathtools}
\usepackage[expansion=false]{microtype}
\usepackage{aliascnt}
\usepackage{pgfplots}
% compat is pinned deliberately: without it pgfplots warns and may change
% defaults between installed versions.
\pgfplotsset{compat=1.18}
\usepackage[colorlinks=true,linkcolor=black,citecolor=black,urlcolor=black]{hyperref}
\usepackage[nameinlink,capitalize]{cleveref}
\hypersetup{
  pdftitle={Sharp coefficientwise positivity for a matrix Tur\'an determinant of 0F1},
  pdfauthor={John Shields},
  pdfkeywords={coefficientwise positivity, Turan inequality, modified Bessel function,
    reciprocal-gamma series, hypergeometric function, order derivative,
    absolute monotonicity}
}
\pdfinfo{ /ORCID (https://orcid.org/0009-0004-3882-4571) }

\numberwithin{equation}{section}

\theoremstyle{plain}
\newtheorem{theorem}{Theorem}[section]
\newaliascnt{proposition}{theorem}
\newtheorem{proposition}[proposition]{Proposition}
\aliascntresetthe{proposition}
\newaliascnt{lemma}{theorem}
\newtheorem{lemma}[lemma]{Lemma}
\aliascntresetthe{lemma}
\newaliascnt{corollary}{theorem}
\newtheorem{corollary}[corollary]{Corollary}
\aliascntresetthe{corollary}

\theoremstyle{definition}
\newaliascnt{remark}{theorem}
\newtheorem{remark}[remark]{Remark}
\aliascntresetthe{remark}
\newaliascnt{question}{theorem}
\newtheorem{question}[question]{Question}
\aliascntresetthe{question}

\crefname{theorem}{theorem}{theorems}
\Crefname{theorem}{Theorem}{Theorems}
\crefname{proposition}{proposition}{propositions}
\Crefname{proposition}{Proposition}{Propositions}
\crefname{lemma}{lemma}{lemmas}
\Crefname{lemma}{Lemma}{Lemmas}
\crefname{corollary}{corollary}{corollaries}
\Crefname{corollary}{Corollary}{Corollaries}
\crefname{remark}{remark}{remarks}
\Crefname{remark}{Remark}{Remarks}
\crefname{question}{question}{questions}
\Crefname{question}{Question}{Questions}

\newcommand{\E}{\mathbb E}
\newcommand{\R}{\mathbb R}
\newcommand{\N}{\mathbb N}
\newcommand{\Dfin}{\mathsf D}
\newcommand{\diag}{\operatorname{diag}}
\newcommand{\MD}{\operatorname{MD}}
\newcommand{\dd}{\,\mathrm d}
\newcommand{\Turan}{\mathcal T}
\newcommand{\Thetal}{\Theta_{\lambda}}
\newcommand{\Thetaz}{\Theta_z}

\title[Sharp coefficientwise positivity for a matrix Tur\'an determinant of \({}_0F_1\)]
{Sharp coefficientwise positivity\\ for a matrix Tur\'an determinant of \({}_0F_1\)}

\author{John Shields}
\date{July 28, 2026}
\thanks{Independent researcher, Ginza 2-15-2, KR Ginza II 5F, Chuo-ku, Tokyo, Japan 104-0061.  Corresponding author: \href{mailto:math@johnnyshields.com}{math@johnnyshields.com}.  ORCID: \href{https://orcid.org/0009-0004-3882-4571}{0009-0004-3882-4571}.  This work is archived at \href{https://doi.org/10.5281/zenodo.21438795}{doi:10.5281/zenodo.21438795}.}

\subjclass[2020]{Primary 33C10; Secondary 15A45, 26A48, 26D15}
\keywords{coefficientwise positivity, Tur\'an inequality, modified Bessel function, reciprocal-gamma series, hypergeometric function, order derivative, absolute monotonicity}

\begin{document}

\begin{abstract}
Let \(Z(a,\lambda)={}_0F_1(;a;\lambda)/\Gamma(a)=\sum_{k\ge0}\lambda^k/(k!\,\Gamma(a+k))\) for \(a>0\).  We prove a sharp coefficientwise positivity theorem for a confluent matrix Tur\'an determinant of this reciprocal-gamma normalization.  From parameter and Euler derivatives of \(Z\), we form a one-parameter family of symmetric \(2\times2\) matrices \(\Turan_\kappa(a,\lambda)\) and prove that, for each fixed \(a>0\), every positive-degree Maclaurin coefficient of \(\Delta^{(\kappa)}=\det\Turan_\kappa\) is strictly positive if and only if \(\kappa\ge1\).  Thus \(\kappa=1\) is the exact deformation threshold at every parameter value, and \(\Delta^{(\kappa)}(a,\cdot)\) is strictly absolutely monotone for every \(\kappa\ge1\).  At this endpoint, coefficientwise determinant positivity holds even though the first nonconstant coefficient matrix may be indefinite: a mixed-determinant polarization confines its possible negative contribution to the degree-two coefficient, whose positivity follows from an explicit decomposition.

Via the identity \(I_{a-1}(2\sqrt\lambda)=\lambda^{(a-1)/2}Z(a,\lambda)\), an exact diagonal congruence transfers the endpoint theorem to the modified Bessel function \(I_\nu\), yielding, for every \(\nu>-1\) and \(z>0\), a positive-definite Bessel--Schur matrix coupling order, argument, and mixed order--argument derivatives of \(\log I_\nu\).  For each fixed \(\nu\), neither the deformation threshold \(\kappa=1\) nor the argument-independent correction \(4/\psi_1(\nu+1)\) can be lowered.
\end{abstract}

\maketitle
\section{Introduction}

Tur\'an inequalities for modified Bessel functions are closely connected with
order monotonicity, logarithmic derivatives, and total positivity; see \cite{Baricz2008,Baricz2010,BariczPonnusamy2013,Watson1944}.
A complementary literature proves scalar log-concavity and coefficientwise
Tur\'an inequalities for reciprocal-gamma and hypergeometric series
\cite{KalmykovKarp2013Reciprocal,KalmykovKarp2013Ratios,KalmykovKarp2014,KarpSitnik2010,KarpZhang2024}.
Here we study the normalized series
\[
 Z(a,\lambda)=\sum_{k\ge0}\frac{\lambda^k}{k!\,\Gamma(a+k)}
\]
through a coupled \(2\times2\) matrix involving parameter and Euler derivatives.

The entries of that matrix, \(\Turan_\kappa(a,\lambda)\) of \eqref{eq:Tkappa}, are
normalized logarithmic derivatives of \(Z\) in the parameter direction
\(\partial_a\) and the Euler direction \(\lambda\partial_\lambda\), together with the
mixed derivative; a real parameter \(\kappa\) scales the Euler-direction term
alone.  Its
determinant \(\Delta^{(\kappa)}=\det\Turan_\kappa\) vanishes at \(\lambda=0\), and
the central result, \Cref{thm:coefficientwise}, decides for which \(\kappa\) all of
its remaining Maclaurin coefficients are positive: for each fixed \(a>0\), every
coefficient of positive degree in \(\Delta^{(\kappa)}(a,\cdot)\) is strictly
positive if and only if \(\kappa\ge1\).  Coefficientwise positivity is strictly
stronger than positivity of \(\Delta^{(\kappa)}\) itself, and at the endpoint
\(\kappa=1\) it makes \(\Delta=\Delta^{(1)}(a,\cdot)\) strictly absolutely monotone
on \((0,\infty)\).  Through the identity \(I_{a-1}(2\sqrt\lambda)
=\lambda^{(a-1)/2}Z(a,\lambda)\), the endpoint case yields \Cref{thm:bessel}: for
every \(\nu>-1\) and \(z>0\),
\[
 \bigl(-\partial_\nu^2\log I_\nu\bigr)
 \left(2\bigl(\Thetaz\log I_\nu-\nu\bigr)-\Thetaz^2\log I_\nu
 +\frac4{\psi_1(\nu+1)}\right)
 >\bigl(1+\partial_\nu\Thetaz\log I_\nu\bigr)^2,
\]
where \(\Thetaz=z\partial_z\) and \(\psi_1\) is the trigamma function.  The matrix
behind this inequality, the Bessel--Schur matrix of
\Cref{cor:bessel-matrix}, couples the concavity of \(\log I_\nu\) in the order, the
argument-direction expression above, and the mixed order--argument sensitivity;
\(\Delta\) is the corresponding differential Schur determinant.  By a
\emph{matrix Tur\'an determinant} we mean this confluent Tur\'anian assembled from
derivatives, rather than a determinant of parameter-shifted evaluations (compared
in \Cref{sec:context}).

A single threshold in \(\kappa\) serves every \(a>0\) and is exact at each \(a\)
separately, the two failures being detected respectively in degree one as
\(a\downarrow0\) and at large argument (\Cref{cor:converse-fixed-a}).  On the
Bessel side, \Cref{prop:bessel-sharpness} identifies \(4/\psi_1(\nu+1)\) as the
least argument-independent correction: no smaller constant keeps the displayed
inequality true for all \(z>0\).  Both sharpness statements hold on the
positive-series domain \(a>0\) (equivalently \(\nu>-1\)), the range in which every
coefficient of \(Z\) is positive.

\Cref{sec:coefficients} opens with an asymmetric Chu--Vandermonde convolution,
which yields exact formulas for the
coefficient matrices of \(\Turan_\kappa\).  These formulas expose the structure
of the sharp endpoint \(\kappa=1\): the convolution identity
singles out \(m=1\) as the potentially exceptional positive degree.

In \Cref{sec:gram}, explicit \(\ell^2\) Gram representations prove that every
coefficient matrix of degree \(m\ge2\) is positive definite, whereas \(M_1\) can
be indefinite for small \(a>0\).  \Cref{sec:determinant} then passes, by
polarization, from positivity of the individual coefficient matrices to
positivity of the determinant coefficients.
The potentially indefinite matrix \(M_1\) can contribute negatively
only through its self-pairing, which occurs solely in the coefficient
\(\Delta_2\); an independent explicit decomposition proves that this exceptional
coefficient is nevertheless strictly positive.  This completes the endpoint
theorem.  Comparison of \(\Turan_\kappa\) with \(\Turan_1\) then gives sufficiency
for \(\kappa\ge1\) and the sharpness of the common threshold.

Finally, in \Cref{sec:bessel-reduction,sec:bessel-consequences}, an explicit
diagonal congruence transports the endpoint theorem to a positive-definite
Bessel--Schur matrix and a mixed order--argument inequality for every
\(\nu>-1\).  The same reduction yields the two sharpness statements of
\Cref{prop:bessel-sharpness} and, in the
large-argument limit, the necessity at fixed \(a\).
\Cref{sec:continuation} examines the boundary of the positive-series domain:
although the
leading small-argument term remains positive away from the poles, the continued
determinant already loses coefficientwise positivity on \(-3<\nu<-2\)
(\Cref{prop:negative-coeff-failure}).

\section{Main results}\label{sec:main}

Throughout, \(a>0\), \(\lambda\ge0\), and \(\N_0=\{0,1,2,\ldots\}\).  We use
\begin{equation}\label{eq:Zdef}
 Z(a,\lambda)=\sum_{k=0}^{\infty}\frac{\lambda^k}{k!\,\Gamma(a+k)}
 =\frac{{}_0F_1(;a;\lambda)}{\Gamma(a)}.
\end{equation}
The series in \eqref{eq:Zdef} and all fixed-order \(a\)- and Euler derivatives
used below converge locally uniformly, and extend to entire functions of
\((a,\lambda)\in\mathbb C^2\)
(\Cref{lem:continuation}); the standing assumptions \(a>0\) and \(\lambda\ge0\) are
in force until \Cref{sec:continuation}, where they are relaxed.  Put
\(\Thetal=\lambda\partial_\lambda\) and \(g=\psi_1(a)\), where
\(\psi=\psi^{(0)}\) and \(\psi_1=\psi^{(1)}\) are the digamma and trigamma
functions.  In the Bessel variable we use the distinct notation
\(\Thetaz=z\partial_z\).  Subscripts denote application of the indicated operator,
so \(Z_a=\partial_aZ\), \(Z_{\Thetal}=\Thetal Z\), and
\(Z_{a\Thetal}=\partial_a\Thetal Z\).  For \(\kappa\in\R\), define
\begin{align}
 A&=Z_a^2-ZZ_{aa},\label{eq:Adef}\\
 B&=Z^2+ZZ_{a\Thetal}-Z_aZ_{\Thetal},\label{eq:Bdef}\\
 C_\kappa&=Z^2+g\,\bigl(\kappa ZZ_{\Thetal}-ZZ_{\Thetal\Thetal}+Z_{\Thetal}^2\bigr),
 \label{eq:Ckappa-def}
\end{align}
\begin{equation}\label{eq:Tkappa}
 \Turan_\kappa=
 \begin{pmatrix}A&\sqrt g\,B\\ \sqrt g\,B&C_\kappa\end{pmatrix},
 \qquad
 \Delta^{(\kappa)}=\det\Turan_\kappa=AC_\kappa-gB^2.
\end{equation}
At \(\kappa=1\) write \(C=C_1\), \(\Turan=\Turan_1\), and \(\Delta=\Delta^{(1)}\).

\begin{remark}\label{rem:canonical-origin}
Writing \(L=\log Z\), direct differentiation gives
\[
 \frac{A}{Z^2}=-L_{aa},\qquad
 \frac{B}{Z^2}=1+L_{a\Thetal},\qquad
 \frac{C_\kappa}{gZ^2}
 =\frac1g+\kappa L_{\Thetal}-L_{\Thetal\Thetal}.
\]
Thus the entries are normalized logarithmic derivatives of \(Z\) in the parameter
and Euler directions --- a curvature in the parameter direction, a shifted mixed
derivative, and a combination of first and second Euler derivatives --- with
\(\kappa\) deforming the argument-direction term.
At the endpoint \(\kappa=1\), the normalization is chosen so that the identity
\eqref{eq:I-Z} and a diagonal congruence carry \(\Turan\) to the
order--argument Bessel--Schur matrix of \Cref{cor:bessel-matrix}; see
\Cref{sec:bessel-reduction}.
\end{remark}

A function on \((0,\infty)\) is \emph{strictly absolutely monotone} if it and all of
its derivatives are strictly positive there.

\begin{theorem}[Sharp coefficientwise classification]\label{thm:coefficientwise}
For each \(a>0\),
\begin{equation}\label{eq:Deltakappa-series}
 \Delta^{(\kappa)}(a,\lambda)
 =\sum_{n\ge1}\Delta_n^{(\kappa)}(a)\lambda^n,
 \qquad \Delta_n^{(\kappa)}(a)>0
 \quad\text{for every }n\ge1,
\end{equation}
holds if and only if \(\kappa\ge1\).  The threshold is therefore exact at every
parameter value, and a single threshold serves every \(a>0\).  For \(\kappa<1\) the
linear coefficient is negative for all sufficiently small \(a>0\), while at a
fixed \(a>0\) the first negative coefficient can occur in arbitrarily high degree
as \(\kappa\uparrow1\) (\Cref{rem:first-negative-degree}); the necessity at fixed
\(a\) is proved in \Cref{cor:converse-fixed-a}.  Consequently, for every \(\kappa\ge1\),
\(\Delta^{(\kappa)}(a,\cdot)\) is strictly absolutely monotone on \((0,\infty)\) and
\(\Turan_\kappa(a,\lambda)\succ0\) for \(a,\lambda>0\).
\end{theorem}

That a single threshold serves every \(a>0\) does not assert a positive lower
bound uniform in \(a\).

The series in \eqref{eq:Deltakappa-series} starts at degree one because
\[
 A(a,0)=\frac{g}{\Gamma(a)^2},\qquad
 B(a,0)=C_\kappa(a,0)=\frac1{\Gamma(a)^2},
\]
so \(\Turan_\kappa(a,0)=\Gamma(a)^{-2}
\left(\begin{smallmatrix}g&\sqrt g\\ \sqrt g&1\end{smallmatrix}\right)\) has
rank one and \(\Delta^{(\kappa)}(a,0)=0\).

\begin{remark}\label{rem:rank-one-threshold}
The rank-one constant term is shared by the entire family
\(\Turan_\kappa\), so the sign of the first nonzero determinant coefficient
detects the sharp threshold in \(\kappa\) uniformly in \(a\), although positivity
in every degree for \(\kappa\ge1\) requires the full coefficientwise argument.
\end{remark}

Under the correspondence \(a=\nu+1\), \(z=2\sqrt\lambda\) of the identity
\eqref{eq:I-Z}, order and argument derivatives of \(I_\nu\) correspond to \(a\)- and
Euler derivatives of \(Z\), which exist by the locally uniform convergence
recorded after \eqref{eq:Zdef} and the entirety in \((a,\lambda)\) of
\Cref{lem:continuation}.  For \(\nu>-1\) and \(z>0\) the positive coefficients of
\(Z(a,\lambda)\) show that \(I_\nu(z)>0\), so \(\log I_\nu\) is real and the following
logarithmic derivatives are well defined:
\begin{align}
 G_\nu&=-\partial_\nu^2\log I_\nu,\label{eq:Gnu}\\
 P_\nu&=\partial_\nu(\Thetaz\log I_\nu),\label{eq:Pnu}\\
 H_\nu&=2(\Thetaz\log I_\nu-\nu)-\Thetaz^2\log I_\nu,
 \label{eq:Hnu}
\end{align}
and
\begin{equation}\label{eq:Dnu-def}
 D_\nu=G_\nu\left(H_\nu+\frac4{\psi_1(\nu+1)}\right)-(1+P_\nu)^2.
\end{equation}
The shifts in \(1+P_\nu\) and \(4/\psi_1(\nu+1)\) are exactly those produced by
the diagonal congruence \eqref{eq:bessel-congruence}.  With \(L=\log Z\) as in
\Cref{rem:canonical-origin},
\(P_\nu=1+2L_{a\Thetal}\), so \(1+P_\nu=2(1+L_{a\Thetal})\); as \(z\downarrow0\) one
has \(L_{a\Thetal}\to0\), hence \(P_\nu\to1\) and \(1+P_\nu\to2\).  For
\(\kappa\in\R\), define the deformation
\begin{align}
 H_\nu^{(\kappa)}
 &=2\kappa(\Thetaz\log I_\nu-\nu)-\Thetaz^2\log I_\nu,
 \label{eq:Hnu-kappa}\\
 D_\nu^{(\kappa)}
 &=G_\nu\left(H_\nu^{(\kappa)}+\frac4{\psi_1(\nu+1)}\right)
 -(1+P_\nu)^2.\label{eq:Dnu-kappa-def}
\end{align}
At \(\kappa=1\) these reduce to \(H_\nu\) and \(D_\nu\).  We call this the
\emph{fixed-correction deformation} because the argument-independent correction
\(4/\psi_1(\nu+1)\) is held fixed while the argument-direction term is deformed.

\begin{theorem}[Mixed Bessel--Schur inequality]\label{thm:bessel}
For \(\nu>-1\) and \(z>0\), \(D_\nu(z)>0\), equivalently
\begin{equation}\label{eq:bessel-main}
 G_\nu(z)\left(H_\nu(z)+\frac4{\psi_1(\nu+1)}\right)>(1+P_\nu(z))^2.
\end{equation}
\end{theorem}

\begin{corollary}[Positive Bessel--Schur matrix]\label{cor:bessel-matrix}
For \(\nu>-1\) and \(z>0\),
\[
 \begin{pmatrix}
 G_\nu(z)&1+P_\nu(z)\\
 1+P_\nu(z)&H_\nu(z)+4/\psi_1(\nu+1)
 \end{pmatrix}\succ0.
\]
As \(z\downarrow0\), the matrix-valued function admits a continuous extension whose value is the rank-one matrix
\(\left(\begin{smallmatrix}g&2\\2&4/g\end{smallmatrix}\right)\), where \(g=\psi_1(\nu+1)\).
\end{corollary}

\begin{proposition}[Sharpness of the Bessel correction and deformation]\label{prop:bessel-sharpness}
For each fixed \(\nu>-1\), \(4/\psi_1(\nu+1)\) is the least real constant \(R\)
such that
\[
 G_\nu(z)\bigl(H_\nu(z)+R\bigr)>(1+P_\nu(z))^2
 \qquad\text{for every }z>0.
\]
Moreover, for the fixed-correction deformation \(D_\nu^{(\kappa)}\) defined in
\eqref{eq:Dnu-kappa-def} and each fixed \(\nu>-1\),
\[
 D_\nu^{(\kappa)}(z)>0
 \qquad\text{for every }z>0
\]
holds if and only if \(\kappa\ge1\).
\end{proposition}

Sufficiency and the sharpness of the common threshold in
\Cref{thm:coefficientwise} are proved in \Cref{sec:determinant}, and the
necessity at each fixed \(a>0\) is completed in \Cref{cor:converse-fixed-a};
\Cref{thm:bessel,cor:bessel-matrix,prop:bessel-sharpness} are proved in
\Cref{sec:bessel-consequences}.


\section{Reciprocal-gamma convolution and coefficient formulas}\label{sec:coefficients}

The coefficient calculation rests on a two-parameter Chu--Vandermonde identity.

\begin{lemma}[Asymmetric reciprocal-gamma convolution]\label{lem:convolution}
Let \(m\in\N_0\) and let \(\alpha,\beta>0\).  Then
\begin{equation}\label{eq:asymmetric-convolution}
 \sum_{i=0}^m
 \frac1{i!(m-i)!\Gamma(\alpha+i)\Gamma(\beta+m-i)}
 =\frac{(\alpha+\beta+m-1)_m}
 {m!\Gamma(\alpha+m)\Gamma(\beta+m)}.
\end{equation}
Both sides are entire in \((\alpha,\beta)\).
\end{lemma}

\begin{proof}
Multiply the left-hand side by \(m!\Gamma(\alpha+m)\Gamma(\beta+m)\).  Using
\[
 \frac{\Gamma(\alpha+m)}{\Gamma(\alpha+i)}=(\alpha+i)_{m-i}
 =\frac{(\alpha)_m}{(\alpha)_i}
\]
and
\[
 \frac{\Gamma(\beta+m)}{\Gamma(\beta+m-i)}
 =(\beta+m-i)_i=(-1)^i(1-\beta-m)_i,
\]
the resulting sum is
\[
 (\alpha)_m
 {}_2F_1(-m,1-\beta-m;\alpha;1).
\]
Chu--Vandermonde gives \((\alpha+\beta+m-1)_m\).  Both sides of
\eqref{eq:asymmetric-convolution} are entire in \((\alpha,\beta)\), being finite
sums and products of the entire function \(1/\Gamma\) and a polynomial, so
agreement on \(\{\alpha,\beta>0\}\) extends the identity to \(\mathbb C^2\) by the
identity theorem in each variable separately.
\end{proof}

Set
\[
 S_m=S_m(a)=\frac{(2a+m-1)_m}{m!\Gamma(a+m)^2}>0.
\]
Taking \(\alpha=\beta=a\) in \Cref{lem:convolution} gives
\[
 [\lambda^m]Z(a,\lambda)^2=S_m,
\]
where \([\lambda^m]\) extracts the coefficient of \(\lambda^m\).  Unsubscripted
\(\alpha,\beta\) always denote the parameters of \Cref{lem:convolution}; the
subscripted families of \Cref{thm:coefficients} below are coefficient sequences.

\begin{theorem}\label{thm:coefficients}
There are expansions
\begin{equation}\label{eq:ABC-expansions}
 A=\sum_{m\ge0}S_m\alpha_m\lambda^m,
 \qquad
 B=\sum_{m\ge0}S_m\beta_m\lambda^m,
 \qquad
 C_\kappa=\sum_{m\ge0}S_m\gamma_m^{(\kappa)}\lambda^m,
\end{equation}
where
\begin{equation}\label{eq:alpha}
 \alpha_m=\psi_1(a+m),
\end{equation}
\begin{equation}\label{eq:beta}
 \beta_0=1,
 \qquad
 \beta_m=\frac{2a+m-2}{2(a+m-1)}\quad(m\ge1),
\end{equation}
and
\begin{equation}\label{eq:gamma-kappa}
 \gamma_m^{(\kappa)}=1+g c_m^{(\kappa)},
\end{equation}
with
\begin{equation}\label{eq:cm-kappa}
 c_0^{(\kappa)}=0,
 \qquad
 c_1^{(\kappa)}=\frac{\kappa-1}{2},
\end{equation}
and, for \(m\ge2\),
\begin{equation}\label{eq:cm-kappa-general}
 c_m^{(\kappa)}
 =\frac{\kappa m}{2}
 -\frac{m(2a+m-2)}{2(2a+2m-3)}
 =\frac{m\left[(\kappa-1)(2a+2m-3)+m-1\right]}
 {2(2a+2m-3)}.
\end{equation}
At the sharp endpoint \(\kappa=1\),
\[
 c_0=c_1=0,
 \qquad
 c_m=\frac{m(m-1)}{2(2a+2m-3)}\quad(m\ge2).
\]
Consequently,
\begin{equation}\label{eq:matrix-series}
 \Turan(a,\lambda)=\sum_{m\ge0}S_mM_m\lambda^m,
 \qquad
 M_m=
 \begin{pmatrix}
 \alpha_m&\sqrt g\,\beta_m\\
 \sqrt g\,\beta_m&1+gc_m
 \end{pmatrix}.
\end{equation}
\end{theorem}

\begin{proof}
For the first coefficient, define
\[
 F_m(\delta)=
 \sum_{i=0}^m
 \frac1{i!(m-i)!\Gamma(a+\delta+i)\Gamma(a-\delta+m-i)}.
\]
By \Cref{lem:convolution},
\begin{equation}\label{eq:Fdelta}
 F_m(\delta)=
 \frac{(2a+m-1)_m}
 {m!\Gamma(a+m+\delta)\Gamma(a+m-\delta)}.
\end{equation}
On the one hand,
\begin{equation}\label{eq:A-second-delta}
 \left.\partial_\delta^2
 \bigl[Z(a+\delta,\lambda)Z(a-\delta,\lambda)\bigr]
 \right|_{\delta=0}
 =2ZZ_{aa}-2Z_a^2=-2A.
\end{equation}
On the other hand, \(F_m'(0)=0\) and the logarithmic second derivative of
\eqref{eq:Fdelta} is \(-2\psi_1(a+m)\) at \(\delta=0\), so
\begin{equation}\label{eq:F-second-delta}
 F_m''(0)=-2S_m\psi_1(a+m).
\end{equation}
Comparing the coefficient of \(\lambda^m\) in
\eqref{eq:A-second-delta} with \eqref{eq:F-second-delta} gives
\([\lambda^m]A=S_m\psi_1(a+m)\), proving \eqref{eq:alpha}.

For \(m=0\), one has \([\lambda^0]B=Z(a,0)^2=S_0\), so \(\beta_0=1\).
Assume henceforth that \(m\ge1\).  Let \(\mathsf I\) be the index of the first factor in
\(Z(a+\delta,\lambda)Z(a-\delta,\lambda)=\sum_mF_m(\delta)\lambda^m\), so that
\([\lambda^m]\) weights \(i\) by
\[
 w_i(\delta)=\bigl[i!(m-i)!\Gamma(a+\delta+i)\Gamma(a-\delta+m-i)\bigr]^{-1}.
\]
For real \(\delta\) with \(|\delta|<a\), write \(\E_\delta\) for expectation
under the probability weights \(w_i(\delta)/F_m(\delta)\), and put
\(\E=\E_0\), \(\Dfin=m-2\mathsf I\), and \(\mathsf J=m-\mathsf I\).  Applying \Cref{lem:convolution} to
\(\sum_i i\,w_i(\delta)\), with \((\alpha,\beta)=(a+1+\delta,a-\delta)\) and \(m-1\) in
place of \(m\), gives
\[
 \sum_{i=0}^m i\,w_i(\delta)
 =\frac{(2a+m-1)_{m-1}}
 {(m-1)!\,\Gamma(a+m+\delta)\Gamma(a+m-1-\delta)}.
\]
Dividing by \eqref{eq:Fdelta} yields
\[
 \E_\delta \mathsf I=\frac{m(a+m-1-\delta)}{2(a+m-1)},
\]
whence
\begin{equation}\label{eq:EDdelta}
 \E_\delta\Dfin=\frac{m\delta}{a+m-1}.
\end{equation}
The score \(\partial_\delta\log w_i(\delta)\big|_{\delta=0}\) is
\[
 X=\psi(a+m-\mathsf I)-\psi(a+\mathsf I),
\]
so differentiating \eqref{eq:EDdelta} at \(\delta=0\) gives
\(\E(\Dfin X)=\partial_\delta\E_\delta\Dfin|_{\delta=0}\), that is
\begin{equation}\label{eq:EDX}
 \E(\Dfin X)=\frac{m}{a+m-1}.
\end{equation}
Direct coefficient extraction gives
\begin{align*}
 &[\lambda^m]\bigl(ZZ_{a\Thetal}-Z_aZ_{\Thetal}\bigr)\\
 &\quad=\sum_{i=0}^m
 \frac{-(m-i)\psi(a+m-i)+(m-i)\psi(a+i)}
 {i!\,(m-i)!\,\Gamma(a+i)\Gamma(a+m-i)}
 =-S_m\E(\mathsf JX).
\end{align*}
Together with \([\lambda^m]Z^2=S_m\), this yields
\[
 \frac{[\lambda^m]B}{S_m}=1-\E(\mathsf JX)
 =1-\frac12\E(\Dfin X),
\]
because \(\mathsf J=(m+\Dfin)/2\) and \(\E X=0\) by symmetry.  This proves
\eqref{eq:beta}.

For \(C_\kappa\), the Euler-derivative terms contribute
\[
 \kappa\E \mathsf J-\E \mathsf J^2+\E(\mathsf I\mathsf J).
\]
Since \(\mathsf I=(m-\Dfin)/2\), \(\mathsf J=(m+\Dfin)/2\), and \(\E\Dfin=0\) by the symmetry \(\mathsf I\leftrightarrow m-\mathsf I\),
\begin{equation}\label{eq:c-cond-derivation}
 \kappa\E \mathsf J-\E \mathsf J^2+\E(\mathsf I\mathsf J)
 =\frac{\kappa m}{2}-\frac12\E\Dfin^2.
\end{equation}
For \(m\ge2\), cancelling the factors \(i(m-i)\) and shifting
\(j=i-1\) gives
\begin{align*}
 S_m(a)\E[\mathsf I(m-\mathsf I)]
 &=\sum_{i=1}^{m-1}
 \frac1{(i-1)!(m-i-1)!\Gamma(a+i)\Gamma(a+m-i)}\\
 &=S_{m-2}(a+1).
\end{align*}
Therefore
\[
 \E[\mathsf I(m-\mathsf I)]
 =\frac{S_{m-2}(a+1)}{S_m(a)}
 =\frac{m(m-1)(a+m-1)}{2(2a+2m-3)},
\]
so
\[
 \E\Dfin^2=\frac{m(2a+m-2)}{2a+2m-3}.
\]
Substitution gives \eqref{eq:cm-kappa-general}; the cases \(m=0,1\) are immediate.
\end{proof}

\begin{remark}\label{rem:coefficient-caveat}
\Cref{thm:coefficients} does not assert \(M_m\succeq0\) for every \(m\): \(M_1\) can be indefinite for small \(a>0\) (see \Cref{sec:gram}).  For general \(\kappa\), write \(M_m^{(\kappa)}\) for the coefficient matrix of \(\Turan_\kappa\), obtained by replacing \(1+gc_m\) with \(1+gc_m^{(\kappa)}\).
\end{remark}

\begin{remark}\label{rem:finite-law}
The weights
\(w_i(0)\) of the proof of \Cref{thm:coefficients}, normalized by \(F_m(0)=S_m\),
form a symmetric finite law on \(0\le i\le m\), under which \(\mathsf I\) is
the random index; \(\E\), \(\Dfin\), and \(X\) are as defined there, and
\[
 \sigma=\psi_1(a+\mathsf I)+\psi_1(a+m-\mathsf I).
\]
The three normalized coefficient entries may be written as
\[
 \alpha_m=\tfrac12\E(\sigma-X^2),\qquad
 \beta_m=1-\tfrac12\E(\Dfin X),\qquad
 c_m^{(\kappa)}=\tfrac{\kappa m}{2}-\tfrac12\E\Dfin^2.
\]
The first identity follows by writing
\[
 \frac{F_m''(0)}{F_m(0)}
 =\E(X^2-\sigma)
 =-2\psi_1(a+m),
\]
while the other two are the expectation identities obtained in \eqref{eq:EDX}
and \eqref{eq:c-cond-derivation}.  Thus the asymmetric Chu--Vandermonde
calculation evaluates, respectively, a curvature term, an imbalance--score
covariance, and an imbalance variance.
\end{remark}

\section{Gram structure and the exceptional matrix \texorpdfstring{\(M_1\)}{M1}}\label{sec:gram}

We require elementary two-sided trigamma estimates.  The series
\(\psi_1(y)=\sum_{r\ge0}(y+r)^{-2}\) is standard
\cite[eq.~5.15.1]{NISTDLMFGamma}; inserting
\((y+r)^{-2}=\int_0^\infty\tau e^{-(y+r)\tau}\dd\tau\) and summing
\(\sum_{r\ge0}e^{-r\tau}=(1-e^{-\tau})^{-1}\) under the integral gives the
integral representation
\begin{equation}\label{eq:trigamma-integral}
 \psi_1(y)=\int_0^\infty e^{-y\tau}\frac{\tau}{1-e^{-\tau}}\dd\tau,
 \qquad y>0.
\end{equation}

\begin{lemma}\label{lem:trigamma-bounds}
For every \(y>0\),
\begin{equation}\label{eq:trig-lower}
 \psi_1(y)>\frac1y+\frac1{2y^2}>\frac1y.
\end{equation}
For \(y>1/2\),
\begin{equation}\label{eq:trig-upper-half}
 \psi_1(y)<\frac1{y-1/2}.
\end{equation}
Consequently,
\begin{equation}\label{eq:inverse-trig}
 \frac1{\psi_1(a)}>a-\frac12
 \qquad(a>0),
\end{equation}
where the assertion is immediate for \(a\le1/2\).
\end{lemma}

\begin{proof}
For \(h(s)=(y+s)^{-2}\), strict convexity on every interval \([r,r+1]\)
gives the strict trapezoidal inequality
\[
 \int_r^{r+1}h(s)\dd s<\frac{h(r)+h(r+1)}2.
\]
Summing and passing to the limit yields
\[
 \psi_1(y)=\sum_{r\ge0}h(r)>
 \int_0^\infty h(s)\dd s+\frac12h(0)
 =\frac1y+\frac1{2y^2},
\]
which proves \eqref{eq:trig-lower}.  For the upper bound, the elementary
inequality \(\tau<2\sinh(\tau/2)\) for \(\tau>0\) is equivalent to
\(\tau/(1-e^{-\tau})<e^{\tau/2}\).  Substitution in \eqref{eq:trigamma-integral} gives
\eqref{eq:trig-upper-half}.  Taking reciprocals proves
\eqref{eq:inverse-trig} when \(a>1/2\).
\end{proof}

Fix \(m\ge1\) and write
\[
 x=a+m,
 \qquad
 q=a+\frac m2-1.
\]
In \(\ell^2(\N_0)\) define
\[
 u_r^{(m)}=\frac1{x+r},
 \qquad
 v_r^{(m)}=\frac{q}{x+r-1},
 \qquad r\ge0.
\]
Then
\begin{equation}\label{eq:u-norm}
 \|u^{(m)}\|^2=\psi_1(x)=\alpha_m,
\end{equation}
\begin{equation}\label{eq:uv-cross}
 \langle u^{(m)},v^{(m)}\rangle
 =q\sum_{r\ge0}\frac1{(x+r-1)(x+r)}
 =\frac q{x-1}=\beta_m,
\end{equation}
and
\begin{equation}\label{eq:v-norm}
 \|v^{(m)}\|^2=q^2\psi_1(x-1).
\end{equation}
Define
\begin{equation}\label{eq:rho-m}
 \rho_m(a)=\frac1g+c_m-q^2\psi_1(x-1).
\end{equation}

For vectors \(\xi,\eta\) in a Hilbert space, \(\operatorname{Gram}(\xi,\eta)\) denotes the symmetric \(2\times2\) matrix with entries \(\langle\xi,\xi\rangle\), \(\langle\xi,\eta\rangle\), \(\langle\eta,\eta\rangle\).  Write \(N_m=\diag(1,g^{-1/2})M_m\diag(1,g^{-1/2})\) for the normalized coefficient matrix.

\begin{theorem}\label{thm:gram}
If either \(m\ge2\) and \(a>0\), or \(m=1\) and \(a\ge1/2\), then \(\rho_m(a)>0\).  In \(\ell^2(\N_0)\oplus\R\), put
\[
 \xi_m=(u^{(m)},0),
 \qquad
 \eta_m=(v^{(m)},\sqrt{\rho_m(a)}).
\]
Then
\begin{align}
 N_m&=\operatorname{Gram}(\xi_m,\eta_m),\label{eq:Nm-gram}\\
 M_m&=\operatorname{Gram}(\xi_m,\sqrt g\,\eta_m).\label{eq:Mm-gram}
\end{align}
In particular, \(M_m\) is positive definite in these ranges.
\end{theorem}

\begin{proof}
Assume first that \(m\ge2\).  Then \(x-1>1/2\), so \Cref{lem:trigamma-bounds} gives
\[
 \rho_m(a)>
 a-\frac12+c_m-\frac{q^2}{x-3/2}.
\]
The right-hand side simplifies exactly to
\[
 a-\frac12+\frac{m(m-1)}{2(2a+2m-3)}
 -\frac{(a+m/2-1)^2}{a+m-3/2}
 =\frac{m-1}{2(2a+2m-3)}>0.
\]

For \(m=1\), \(q=a-1/2\) and \(c_1=0\).  If \(a>1/2\), then
\[
 q^2\psi_1(a)<q<\frac1{\psi_1(a)},
\]
by \eqref{eq:trig-upper-half} and \eqref{eq:inverse-trig}.  Hence \(\rho_1(a)>0\).  At \(a=1/2\), \(q=0\) and the conclusion is immediate.  The Gram identities now follow from \eqref{eq:u-norm}--\eqref{eq:rho-m}.
In a vanishing combination \(\mu\,\xi_m+\omega\sqrt g\,\eta_m=0\) the last coordinate
gives \(\omega\sqrt{g\rho_m(a)}=0\), hence \(\omega=0\), and then \(\|u^{(m)}\|^2=\psi_1(x)>0\)
forces \(\mu=0\).  The two Gram vectors of \eqref{eq:Mm-gram} are therefore linearly
independent, so \(M_m\succ0\).
\end{proof}

The constant matrix is
\[
 M_0=
 \begin{pmatrix}g&\sqrt g\\ \sqrt g&1\end{pmatrix}
 =\begin{pmatrix}\sqrt g\\1\end{pmatrix}
 \begin{pmatrix}\sqrt g&1\end{pmatrix},
\]
a rank-one positive-semidefinite matrix.  With \(\psi_1(a+1)=g-a^{-2}\) and \(\beta_1=(2a-1)/(2a)\), the remaining coefficient matrix \(M_1\) has determinant
\begin{equation}\label{eq:det-M1}
 \det M_1
 =\psi_1(a+1)-g\beta_1^2
 =\frac{(4a-1)g-4}{4a^2}.
\end{equation}

\begin{lemma}\label{lem:M1-indefinite}
There is \(a_*=0.3690738484\ldots\in(0,1/2)\) such that \(M_1\) is indefinite for
\(0<a<a_*\), positive semidefinite of rank one at \(a=a_*\), and positive definite
for \(a>a_*\).
\end{lemma}

\begin{proof}
Set \(f(a)=(4a-1)\psi_1(a)-4\), so that \(\det M_1\) has the sign of \(f(a)\) by \eqref{eq:det-M1}.  Since \(g=\psi_1(a)=\sum_{r\ge0}(a+r)^{-2}\), for \(0<a<1/2\)
\[
 f'(a)=2\sum_{r\ge0}\frac{2r+1-2a}{(a+r)^3}>0,
\]
while \(f(a)\to-\infty\) as \(a\downarrow0\) and \(f(1/2)=\psi_1(1/2)-4=\pi^2/2-4>0\).  Hence \(f\) has a unique zero \(a_*\in(0,1/2)\).  The trichotomy follows with \Cref{thm:gram} for \(a\ge1/2\) and the positive diagonal of \(M_1\).
\end{proof}

Among the coefficient matrices only \(M_1\) can be indefinite at the endpoint
\(\kappa=1\), so the determinant positivity of \Cref{thm:coefficientwise} cannot
follow from pointwise semidefiniteness of the \(M_m\) alone; see
\Cref{fig:det-M1}.

\begin{figure}[ht]
\centering
\begin{tikzpicture}
\begin{axis}[
    width=9.4cm, height=6.4cm,
    xlabel={\(a\)}, ylabel={\(f(a)\)},
    xmin=0.25, xmax=0.5, ymin=-4.35, ymax=1.6,
    xtick={0.25,0.3,0.35,0.4,0.45,0.5},
    xticklabel style={/pgf/number format/fixed,
                      /pgf/number format/fixed zerofill,
                      /pgf/number format/precision=2},
    ytick={-4,-3,-2,-1,0,1},
    tick label style={font=\footnotesize},
    label style={font=\small},
    axis lines=left,
    axis on top,
]
\addplot[black!35, dashed, line width=0.6pt, forget plot]
  coordinates {(0.25,0) (0.5,0)};
\addplot[black!35, dashed, line width=0.6pt, forget plot]
  coordinates {(0.36907385,-4.35) (0.36907385,1.6)};
\addplot[black!70, line width=1.0pt] coordinates {(0.25000000,-4.00000000) (0.25250000,-3.83121245) (0.25500000,-3.66861149) (0.25750000,-3.51193050) (0.26000000,-3.36091644) (0.26250000,-3.21532906) (0.26500000,-3.07494013) (0.26750000,-2.93953276) (0.27000000,-2.80890072) (0.27250000,-2.68284788) (0.27500000,-2.56118757) (0.27750000,-2.44374211) (0.28000000,-2.33034228) (0.28250000,-2.22082684) (0.28500000,-2.11504211) (0.28750000,-2.01284159) (0.29000000,-1.91408549) (0.29250000,-1.81864045) (0.29500000,-1.72637918) (0.29750000,-1.63718010) (0.30000000,-1.55092709) (0.30250000,-1.46750919) (0.30500000,-1.38682032) (0.30750000,-1.30875907) (0.31000000,-1.23322840) (0.31250000,-1.16013550) (0.31500000,-1.08939150) (0.31750000,-1.02091133) (0.32000000,-0.954613509) (0.32250000,-0.890419983) (0.32500000,-0.828255943) (0.32750000,-0.768049679) (0.33000000,-0.709732429) (0.33250000,-0.653238241) (0.33500000,-0.598503838) (0.33750000,-0.545468494) (0.34000000,-0.494073915) (0.34250000,-0.444264129) (0.34500000,-0.395985374) (0.34750000,-0.349186002) (0.35000000,-0.303816383) (0.35250000,-0.259828812) (0.35500000,-0.217177423) (0.35750000,-0.175818109) (0.36000000,-0.135708446) (0.36250000,-0.0968076125) (0.36500000,-0.0590763276) (0.36750000,-0.0224767792) (0.37000000,0.0130274373) (0.37250000,0.0474713798) (0.37500000,0.0808888164) (0.37750000,0.113312280) (0.38000000,0.144773120) (0.38250000,0.175301553) (0.38500000,0.204926708) (0.38750000,0.233676671) (0.39000000,0.261578531) (0.39250000,0.288658417) (0.39500000,0.314941539) (0.39750000,0.340452225) (0.40000000,0.365213954) (0.40250000,0.389249394) (0.40500000,0.412580431) (0.40750000,0.435228200) (0.41000000,0.457213116) (0.41250000,0.478554901) (0.41500000,0.499272611) (0.41750000,0.519384663) (0.42000000,0.538908856) (0.42250000,0.557862399) (0.42500000,0.576261931) (0.42750000,0.594123542) (0.43000000,0.611462795) (0.43250000,0.628294746) (0.43500000,0.644633961) (0.43750000,0.660494535) (0.44000000,0.675890111) (0.44250000,0.690833895) (0.44500000,0.705338669) (0.44750000,0.719416813) (0.45000000,0.733080315) (0.45250000,0.746340785) (0.45500000,0.759209470) (0.45750000,0.771697268) (0.46000000,0.783814735) (0.46250000,0.795572106) (0.46500000,0.806979296) (0.46750000,0.818045920) (0.47000000,0.828781297) (0.47250000,0.839194465) (0.47500000,0.849294188) (0.47750000,0.859088966) (0.48000000,0.868587042) (0.48250000,0.877796416) (0.48500000,0.886724848) (0.48750000,0.895379868) (0.49000000,0.903768785) (0.49250000,0.911898690) (0.49500000,0.919776471) (0.49750000,0.927408811) (0.50000000,0.934802201)};
\addplot[black, mark=*, only marks, mark size=1.7pt, forget plot]
  coordinates {(0.36907385,0)};
\node[font=\footnotesize, anchor=north west, align=left]
  at (axis cs:0.3775,-0.24)
  {\(a_*=0.3690738484\ldots\)\\[-1pt] \(\det M_1=0\), \(M_1\succeq0\)};
\node[font=\footnotesize, anchor=north]
  at (axis cs:0.3085,1.52) {\(M_1\) indefinite};
\node[font=\footnotesize, anchor=north]
  at (axis cs:0.4345,1.52) {\(M_1\succ0\)};
\end{axis}
\end{tikzpicture}
\caption{By \eqref{eq:det-M1}, \(f(a)=(4a-1)\psi_1(a)-4=4a^2\det M_1\).  The
plotted window is fixed by the exact values \(f(1/4)=-4\) and
\(f(1/2)=\pi^2/2-4\).  Strict increase on \((0,1/2)\) makes the zero
\(a_*=0.3690738484\ldots\) unique; \(M_1\) is indefinite for \(a<a_*\), singular
positive semidefinite at \(a=a_*\), and positive definite for \(a>a_*\)
(\Cref{lem:M1-indefinite}).  Nevertheless, at \(\kappa=1\), every determinant
coefficient \(\Delta_n(a)\), \(n\ge1\), is strictly positive for all \(a>0\)
(\Cref{thm:coefficientwise}).}
\label{fig:det-M1}
\end{figure}

\section{Mixed determinants and coefficientwise positivity}\label{sec:determinant}

For symmetric \(2\times2\) matrices \(X=(X_{ij})\) and \(Y=(Y_{ij})\), define
\[
 \MD(X,Y)=X_{11}Y_{22}+X_{22}Y_{11}-2X_{12}Y_{12}.
\]
This is the polarization of the determinant.

\begin{lemma}\label{lem:MD-positive}
If \(X,Y\succeq0\), then \(\MD(X,Y)\ge0\).  If in addition \(X\ne0\) and \(Y\succ0\), then \(\MD(X,Y)>0\).
\end{lemma}

\begin{proof}
Write \(X=\sum_r x_rx_r^T\) and \(Y=\sum_s y_sy_s^T\).  Then
\[
 \MD(X,Y)=\sum_{r,s}
 \left(x_{r,1}y_{s,2}-x_{r,2}y_{s,1}\right)^2.
\]
If this sum vanished and \(X\ne0\), fix a nonzero Gram vector \(x_r\).  Every Gram vector \(y_s\) would then be parallel to \(x_r\), forcing \(Y\) to have rank at most one, contrary to \(Y\succ0\).
\end{proof}

Expanding the determinant of \eqref{eq:matrix-series} gives
\begin{equation}\label{eq:Delta-n-MD}
 \Delta_n
 =\frac12\sum_{k=0}^nS_kS_{n-k}\MD(M_k,M_{n-k}).
\end{equation}
By \Cref{thm:coefficients} and \Cref{rem:coefficient-caveat},
\(\Turan_\kappa(a,\lambda)=\sum_{m\ge0}S_mM_m^{(\kappa)}\lambda^m\), so
\eqref{eq:Delta-n-MD} holds for \(\Delta_n^{(\kappa)}\) with each \(M_m\) replaced by
\(M_m^{(\kappa)}\).

\begin{lemma}\label{lem:Delta-n-noneq2}
For every \(a>0\) and every \(n\ge1\) with \(n\ne2\), one has \(\Delta_n(a)>0\).
\end{lemma}

\begin{proof}
The linear coefficient is strictly positive: \eqref{eq:trig-lower} gives
\(\psi_1(a)>1/a\), so \(ag>1\) and
\[
 \MD(M_0,M_1)=\frac{ag-1}{a^2}>0,
\]
whence
\begin{equation}\label{eq:Delta1-sharp}
 \Delta_1=S_0S_1\MD(M_0,M_1)
 =\frac{2(ag-1)}{a^3\Gamma(a)^4}>0.
\end{equation}
For \(m\ge2\) the mixed term satisfies
\begin{equation}\label{eq:M1-Mm-positive}
 \MD(M_1,M_m)
 =\alpha_1(1+gc_m)+\alpha_m-2g\beta_1\beta_m>0.
\end{equation}
If \(a\ge1/2\), then \(M_1,M_m\succ0\) by \Cref{thm:gram} and the strict case of
\Cref{lem:MD-positive} applies; if \(0<a<1/2\), then \(\beta_1<0<\beta_m\) by
\eqref{eq:beta}, and since every diagonal entry is positive the right-hand side
is strictly positive.
For \(n\ge3\) the self-pair \((1,1)\) does not occur.  Classify the pairs \((k,n-k)\)
in \eqref{eq:Delta-n-MD}: if neither index is \(1\), both matrices in the pair are
positive semidefinite---\(M_0\) is rank-one positive semidefinite and \(M_m\succ0\)
for \(m\ge2\) by \Cref{thm:gram}---so the nonstrict case of \Cref{lem:MD-positive}
gives a nonnegative summand; if one index is \(1\), the other is at least \(2\) and
\eqref{eq:M1-Mm-positive} applies.  Every summand is therefore nonnegative, and
the pair \((M_0,M_n)\) is strictly positive because \(M_0\ne0\) and \(M_n\succ0\).
Hence \(\Delta_n>0\).
\end{proof}

The self-pairing \(\MD(M_1,M_1)=2\det M_1\) occurs solely in \(\Delta_2\), so the
indefiniteness established in \Cref{lem:M1-indefinite} can contribute negatively
only there.

\subsection{The exceptional determinant coefficient \texorpdfstring{\(\Delta_2\)}{Delta2}}\label{subsec:delta2}

\begin{lemma}\label{lem:Delta2-positive}
For every \(a>0\), one has \(\Delta_2(a)>0\).  More precisely, let
\[
 t=\psi_1(a+1),
 \qquad
 \mathcal R_a=t-\frac1{a+1}.
\]
Then
\begin{equation}\label{eq:Delta2-Q}
 \Delta_2=
 \frac{Q_*(a,t)}{2a^6(a+1)^3\Gamma(a)^4},
\end{equation}
where
\begin{align}
 Q_*(a,t)={}&2a^4(a+1)^2\mathcal R_a^2\notag\\
 &+2a^2(a+1)^2(8a^2+3a+1)\mathcal R_a\notag\\
 &+2a(a+1)(5a+3).\label{eq:Qstar-decomp}
\end{align}
\end{lemma}

\begin{proof}
From \eqref{eq:Delta-n-MD},
\begin{equation}\label{eq:Delta2-unsimplified}
 \Delta_2=S_0S_2\MD(M_0,M_2)+S_1^2\det M_1.
\end{equation}
Here
\[
 S_0=\frac1{\Gamma(a)^2},\qquad
 S_1=\frac2{a\Gamma(a)^2},\qquad
 S_2=\frac{2a+1}{a^2(a+1)\Gamma(a)^2},
\]
and, recalling
\(\det M_1=t-g\beta_1^2\) from \eqref{eq:det-M1},
\[
 \MD(M_0,M_2)
 =g+\frac{g^2}{2a+1}+\psi_1(a+2)
   -2g\frac{a}{a+1}.
\]
Substituting \(g=t+a^{-2}\) and \(\psi_1(a+2)=t-(a+1)^{-2}\) into
\eqref{eq:Delta2-unsimplified} and collecting over the denominator in
\eqref{eq:Delta2-Q} first gives
\begin{align*}
 Q_*(a,t)={}&2a^4(a+1)^2t^2\\
 &+2a^2(8a^4+17a^3+13a^2+5a+1)t\\
 &+2a(-8a^4-10a^3+a^2+7a+3).
\end{align*}
After multiplication by the common denominator in \eqref{eq:Delta2-Q}, the
preceding formula is a polynomial identity in \((a,t)\).  Replacing \(t\) by
\(\mathcal R_a+(a+1)^{-1}\) and collecting this polynomial yields the identity
\eqref{eq:Qstar-decomp}.

Applying \eqref{eq:trig-lower} at \(y=a+1\) gives
\(\mathcal R_a=\psi_1(a+1)-1/(a+1)>1/(2(a+1)^2)>0\).
All three coefficients in \eqref{eq:Qstar-decomp} are positive for \(a>0\),
so \(\Delta_2>0\).
\end{proof}

\begin{proof}[Proof of \Cref{thm:coefficientwise}, except fixed-\(a\) necessity]
For every \(a>0\) and \(n\ge1\), the endpoint assertion \(\Delta_n^{(1)}(a)>0\) holds
by \Cref{lem:Delta-n-noneq2} when \(n\ne2\) and by \Cref{lem:Delta2-positive} when
\(n=2\).  Since \(A\) and \(Z\) have positive coefficients and
\[
 [\lambda^n]Z_{\Thetal}=\frac{n}{n!\Gamma(a+n)}>0\qquad(n\ge1),
\]
every positive-degree coefficient of \(AZZ_{\Thetal}\) is strictly positive: in
degree \(n\), for example, the product of the constant coefficients of \(A\) and
\(Z\) with \([\lambda^n]Z_{\Thetal}\) is positive.  Therefore
\[
 \Delta^{(\kappa)}
 =\Delta^{(1)}+(\kappa-1)gA\,ZZ_{\Thetal}
\]
shows that all coefficients remain strictly positive for \(\kappa>1\).

For the converse, the general \(m=1\) coefficient gives
\begin{equation}\label{eq:MD01-kappa}
 \MD(M_0,M_1^{(\kappa)})
 =\frac{\frac{\kappa-1}{2}(ag)^2+ag-1}{a^2}.
\end{equation}
If \(\kappa<1\), then \(ag\to+\infty\) as \(a\downarrow0\), so \eqref{eq:MD01-kappa}, and hence \(\Delta_1^{(\kappa)}\), is negative for all sufficiently small \(a>0\).  This establishes sufficiency for every \(a>0\), and necessity in the form that no \(\kappa<1\) serves every \(a>0\); necessity at each fixed \(a>0\) is proved in \Cref{cor:converse-fixed-a}.

Fix \(\kappa\ge1\).  Since \(\Delta^{(\kappa)}(a,\cdot)\) is entire, each differentiated Maclaurin series converges locally uniformly, so for every \(r\ge0\) and \(\lambda>0\),
\[
 \partial_\lambda^r\Delta^{(\kappa)}(a,\lambda)
 =\sum_{n\ge\max\{1,r\}}
 \frac{n!}{(n-r)!}\Delta_n^{(\kappa)}(a)\lambda^{n-r}>0,
\]
which makes \(\Delta^{(\kappa)}(a,\cdot)\) strictly absolutely monotone.  The leading principal entry satisfies \(A>0\) for \(\lambda>0\) by \eqref{eq:ABC-expansions} and \eqref{eq:alpha}; since \(\Delta^{(\kappa)}>0\), it follows that \(\Turan_\kappa\succ0\).
\end{proof}

Combining \eqref{eq:Delta1-sharp} with \eqref{eq:trig-lower}, which gives
\(ag-1>(2a)^{-1}\), yields the explicit lower bound
\[
 \Delta(a,\lambda)>\Delta_1(a)\lambda
 =\frac{2\bigl(a\psi_1(a)-1\bigr)}{a^3\Gamma(a)^4}\,\lambda
 >\frac{\lambda}{a^4\Gamma(a)^4}
 \qquad(a,\lambda>0).
\]

\section{Exact reduction from \texorpdfstring{\({}_0F_1\)}{0F1} to the Bessel inequality}\label{sec:bessel-reduction}

\Cref{sec:main} defines \(G_\nu\), \(P_\nu\), \(H_\nu\), and \(D_\nu\) on the
Bessel side; the congruence below identifies them with the entries of
\(\Turan\) of \eqref{eq:Tkappa}.

The standard identity \cite[eq.~10.25.2]{NISTDLMFBessel},\,\cite[\S3.7]{Watson1944}
\begin{equation}\label{eq:I-Z}
 I_{a-1}(2\sqrt\lambda)=\lambda^{(a-1)/2}Z(a,\lambda)
\end{equation}
with \(L=\log Z\), \(z=2\sqrt\lambda\), \(\nu=a-1\), and \(\Thetaz=2\Thetal\) gives
\begin{align}
 \Thetaz\log I_\nu&=\nu+2L_{\Thetal}, &
 P_\nu&=1+2L_{a\Thetal},\label{eq:U-L}\\
 G_\nu&=-L_{aa}, & H_\nu&=4(L_{\Thetal}-L_{\Thetal\Thetal}).\label{eq:G-L}
\end{align}
Specializing \Cref{rem:canonical-origin} to \(\kappa=1\),
\begin{equation}\label{eq:ABC-log}
 \frac A{Z^2}=-L_{aa},\qquad
 \frac B{Z^2}=1+L_{a\Thetal},\qquad
 \frac C{gZ^2}=\frac1g+L_{\Thetal}-L_{\Thetal\Thetal}.
\end{equation}
The Bessel--Schur matrix is therefore the congruence
\begin{equation}\label{eq:bessel-congruence}
 \begin{pmatrix}
  G_\nu&1+P_\nu\\
  1+P_\nu&H_\nu+4/g
 \end{pmatrix}
 =\frac1{Z^2}
 \begin{pmatrix}1&0\\0&2/\sqrt g\end{pmatrix}
 \Turan(a,\lambda)
 \begin{pmatrix}1&0\\0&2/\sqrt g\end{pmatrix},
\end{equation}
whose determinant is
\begin{equation}\label{eq:D-Delta}
 D_{a-1}(2\sqrt\lambda)=\frac{4\Delta(a,\lambda)}{gZ(a,\lambda)^4}.
\end{equation}
This reduces \Cref{thm:bessel} and \Cref{cor:bessel-matrix} to the endpoint \(\kappa=1\) of \Cref{thm:coefficientwise}.

\section{Bessel consequences and sharpness}\label{sec:bessel-consequences}

The endpoint coefficientwise positivity proved in \Cref{sec:determinant}, transported through the exact reduction of \Cref{sec:bessel-reduction}, gives the sharp Bessel inequality; the boundary limits at \(z\downarrow0\) and \(z\to\infty\) fix its optimal constants and complete \Cref{thm:coefficientwise}.

\begin{proof}[Proof of \Cref{thm:bessel} and \Cref{cor:bessel-matrix}]
The congruence \eqref{eq:bessel-congruence} and the endpoint \(\kappa=1\) of
\Cref{thm:coefficientwise}, proved in \Cref{sec:determinant}, give positive
definiteness for \(z>0\), and
\eqref{eq:D-Delta} gives the determinant inequality.  As \(z\downarrow0\),
equivalently \(\lambda\downarrow0\), the power series for \(L=\log Z\) gives \(L_{\Thetal},L_{a\Thetal},L_{\Thetal\Thetal}\to0\).  Hence \eqref{eq:U-L}--\eqref{eq:G-L} imply
\[
 G_\nu\longrightarrow g,\qquad
 1+P_\nu\longrightarrow2,\qquad
 H_\nu+\frac4g\longrightarrow\frac4g,
\]
which is the matrix limit stated in \Cref{cor:bessel-matrix}.
\end{proof}

\begin{lemma}\label{lem:large-argument-limit}
For every \(\nu>-1\) and every \(\kappa\in\R\),
\begin{equation}\label{eq:D-large-z-limit}
 D_\nu^{(\kappa)}(z)\longrightarrow2(\kappa-1)
 \qquad(z\to\infty),
\end{equation}
a limit independent of \(\nu\).
\end{lemma}

\begin{proof}
Differentiating an asymptotic expansion of \(\log I_\nu\) requires a domain on
which it is holomorphic, and \(I_\nu\) has zeros.  Fix \(\varphi\in(0,\pi/2)\), put
\(\Sigma_\varphi=\{w:|\arg w|\le\pi/2-\varphi\}\), let \(K\subset\mathbb C\) be a
closed disc about the given order, and set
\[
 E_\mu(w)=(2\pi w)^{1/2}e^{-w}I_\mu(w)
 =1-\frac{4\mu^2-1}{8w}+O(w^{-2}),
\]
the expansion holding uniformly for \(\mu\in K\) and \(w\in\Sigma_\varphi\)
\cite[eqs.~10.17.1 and 10.40.1, \S10.40(iii)]{NISTDLMFBessel}.  Hence
there is \(W=W(K,\varphi)\) with \(|E_\mu(w)-1|<1/2\) whenever \(\mu\in K\),
\(w\in\Sigma_\varphi\), and \(|w|\ge W\).  Write \(\Omega\) for the interior of
\(K\times(\Sigma_\varphi\cap\{|w|\ge W\})\).  On \(\Omega\) the values of \(E_\mu\) lie
in the disc \(|\zeta-1|<1/2\), on which the principal logarithm is holomorphic, so
the composition \(\log E_\mu(w)\) is holomorphic in \((\mu,w)\) and vanishes as
\(w\to\infty\); adding
\(w-\tfrac12\log(2\pi w)\) with the principal branch of \(\log w\) gives a
holomorphic \(\log I_\mu(w)\) on \(\Omega\), which agrees with the real \(\log I_\nu(z)\)
for \(\nu\in K\cap(-1,\infty)\) and \(z>0\).  Taking logarithms in the expansion of
\(E_\mu\) gives
\begin{equation}\label{eq:log-E-expansion}
 \log E_\mu(w)=-\frac{4\mu^2-1}{8w}+O(w^{-2}),
\end{equation}
uniformly on \(\Omega\).  Choose
\(\varepsilon=\varepsilon(\varphi)\in(0,1)\) so that \(|w-z|\le\varepsilon z\) forces
\(w\in\Sigma_\varphi\); for \(z\ge2W/(1-\varepsilon)\) such a disc also
lies in \(|w|>W\).  Let \(R(\mu,w)\) denote the \(O(w^{-2})\) remainder in
\eqref{eq:log-E-expansion}.  Only the given order is needed, so the derivatives
are bounded at the centre \(\mu=\nu\) of \(K\): Cauchy's estimate---on the circle
about \(\nu\) of half the radius of \(K\), and on the disc \(|w-z|\le\varepsilon z\),
whose radius is proportional to \(|w|\)---gives
\[
 \partial_\mu^j\partial_w^\ell R(\nu,z)=O(z^{-2-\ell})
 \qquad(0\le j,\ell\le2),
\]
so each \(w\)-derivative gains a factor \(z^{-1}\) while \(\mu\)-derivatives preserve
the order, and \eqref{eq:log-E-expansion} may be differentiated termwise twice
in each variable at \(\mu=\nu\).

At \(w=z>0\), \eqref{eq:log-E-expansion} therefore gives
\[
 \log I_\nu(z)=z-\frac12\log(2\pi z)-\frac{4\nu^2-1}{8z}+O(z^{-2}),
\]
and hence \(G_\nu=z^{-1}+O(z^{-2})\), \(P_\nu=\nu/z+O(z^{-2})\), and
\(H_\nu^{(\kappa)}=(2\kappa-1)z-\kappa(1+2\nu)+O(z^{-1})\), which give
\eqref{eq:D-large-z-limit}.
\end{proof}

\begin{proof}[Proof of \Cref{prop:bessel-sharpness}]
For the least argument-independent correction, replace \(4/\psi_1(\nu+1)\) in
\eqref{eq:bessel-main} by an argument-independent \(R\) and write \(D_{\nu,R}\) for the
resulting determinant.  Then
\[
 D_{\nu,R}(z)\longrightarrow\psi_1(\nu+1)R-4
 \qquad(z\downarrow0),
\]
so validity for every \(z>0\) forces \(R\ge4/\psi_1(\nu+1)\).  Conversely
\(D_{\nu,R}=D_\nu+G_\nu\bigl(R-4/\psi_1(\nu+1)\bigr)>0\) whenever
\(R\ge4/\psi_1(\nu+1)\), since \(D_\nu>0\) by \Cref{thm:bessel} and
\(G_\nu=A/Z^2>0\) for \(\nu>-1\), \(z>0\).

For the fixed-correction deformation
\eqref{eq:Hnu-kappa}--\eqref{eq:Dnu-kappa-def}, the same reduction gives
\begin{equation}\label{eq:Dkappa-Delta}
 D_{a-1}^{(\kappa)}(2\sqrt\lambda)
 =\frac{4\Delta^{(\kappa)}(a,\lambda)}{gZ(a,\lambda)^4}.
\end{equation}
For each fixed
\(\nu>-1\), necessity of \(\kappa\ge1\) follows from
\Cref{lem:large-argument-limit}: the limit \eqref{eq:D-large-z-limit} is negative
when \(\kappa<1\), so \(D_\nu^{(\kappa)}(z)<0\) for all sufficiently large \(z\).
Sufficiency for \(\kappa\ge1\) is the part of \Cref{thm:coefficientwise} proved in
\Cref{sec:determinant}, and transfers through \eqref{eq:Dkappa-Delta}.
\end{proof}

\begin{corollary}\label{cor:converse-fixed-a}
Fix \(a>0\) and \(\kappa<1\).  Then \(\Delta_n^{(\kappa)}(a)<0\) for some \(n\ge1\).
Consequently, at every fixed \(a>0\), coefficientwise positivity in every positive
degree holds if and only if \(\kappa\ge1\), which completes
\Cref{thm:coefficientwise}.
\end{corollary}

\begin{proof}[Proof of \Cref{cor:converse-fixed-a}, completing \Cref{thm:coefficientwise}]
The limit \eqref{eq:D-large-z-limit} of \Cref{lem:large-argument-limit} is
negative for \(\kappa<1\),
so \(D_{a-1}^{(\kappa)}(z)<0\) for all large \(z\).  Hence \eqref{eq:Dkappa-Delta}
gives \(\Delta^{(\kappa)}(a,\lambda)<0\) for all large \(\lambda>0\).  A power series
with nonnegative coefficients is nonnegative on \((0,\infty)\), so some
\(\Delta_n^{(\kappa)}(a)\) is negative.  Sufficiency for \(\kappa\ge1\) is proved in
\Cref{sec:determinant}.
\end{proof}

\begin{remark}\label{rem:first-negative-degree}
The degree-one argument of \Cref{sec:determinant} exhibits failure as
\(a\downarrow0\).  Fix \(a>0\); then no bounded set of degrees detects every
\(\kappa<1\), since the degree of the first negative coefficient tends to infinity
as \(\kappa\uparrow1\).  In
\(\Delta^{(\kappa)}=\Delta^{(1)}+(\kappa-1)gA\,ZZ_{\Thetal}\) both series have
strictly positive coefficients in every positive degree, so for each \(N\) the
finitely many inequalities \(\Delta_n^{(\kappa)}(a)>0\), \(1\le n\le N\), hold once
\(\kappa<1\) is close enough to \(1\).
\end{remark}

\section{Continuation beyond the positive-series domain}\label{sec:continuation}

The behavior changes qualitatively once the order crosses the endpoint
\(\nu=-1\) of the positive-series domain.  On the open set
\[
 \mathcal U=\bigl\{(\nu,z)\in
 (\R\setminus\{-1,-2,\ldots\})\times(0,\infty):I_\nu(z)\ne0\bigr\},
\]
define \(G_\nu\), \(P_\nu\), \(H_\nu\), and \(D_\nu\) by
\eqref{eq:Gnu}--\eqref{eq:Dnu-def} with \(\log I_\nu\) replaced by \(\log|I_\nu(z)|\);
the order derivatives are taken on \(\mathcal U\), and for each fixed \(\nu\) the
resulting formulas restrict to any connected component of
\(\{z>0:I_\nu(z)\ne0\}\).  Where \(I_\nu>0\) this agrees with the ordinary
logarithm.

For every real \(a\notin\{0,-1,-2,\ldots\}\), the partial-fraction expansion
\begin{equation}\label{eq:trigamma-partial-fraction}
 \psi_1(a)=\sum_{r=0}^{\infty}\frac1{(a+r)^2}
\end{equation}
holds and shows that \(g=\psi_1(a)>0\) \cite[eq.~5.15.1]{NISTDLMFGamma}.

\begin{lemma}\label{lem:continuation}
The series \eqref{eq:Zdef} and all of its fixed-order parameter and Euler
derivatives converge locally uniformly on \(\mathbb C^2\), so \(Z\) and those
derivatives are entire in \((a,\lambda)\).
For each fixed coefficient index, the scalar coefficient identities in
\eqref{eq:ABC-expansions}--\eqref{eq:cm-kappa-general}, and every resulting
finite-convolution formula for a coefficient of \(\Delta\), extend
meromorphically in \(a\).  A formula carrying a product \(S_mc_m^{(\kappa)}\) is to
be read as that complete product, whose apparent pole at \(2a+2m-3=0\) is
cancelled by the zero of \(S_m\) there and is therefore removable.  This
assertion concerns the scalar entries and
\(\Delta\); it does not require a single-valued meromorphic choice of
\(\sqrt{\psi_1(a)}\) on the complex plane.  If
\(a\in\R\setminus\{0,-1,-2,\ldots\}\), \(\lambda>0\), and
\(Z(a,\lambda)\ne0\), then \eqref{eq:ABC-log}, \eqref{eq:bessel-congruence}, and
\eqref{eq:D-Delta} remain valid, with the positive real square root of
\(\psi_1(a)\) and with logarithmic derivatives interpreted through \(\log|Z|\)
and \(\log|I_{a-1}|\) on the relevant nonvanishing component.
\end{lemma}

\begin{proof}
Each summand of \eqref{eq:Zdef} is entire in \(a\).  Uniform Stirling estimates
for \(a\) in compact sets, together with Cauchy estimates for the derivatives of
\(1/\Gamma\), give locally uniform convergence of the series and of every
fixed-order parameter and Euler derivative on compact subsets of \(\mathbb C^2\);
hence \(Z\) and these derivatives are entire in \((a,\lambda)\).  For each fixed
coefficient index, both sides of each scalar coefficient identity are
meromorphic functions of the single variable \(a\) and agree on \(a>0\); since \(a>0\)
has a limit point, the one-variable identity theorem shows they agree throughout
their common domain, the complement in \(\mathbb C\) of the discrete pole set.  Each coefficient of
\(\Delta=AC-gB^2\) is a finite convolution of such scalar coefficient functions
and is therefore meromorphic as well.

On the real nonpole set, \eqref{eq:trigamma-partial-fraction} gives
\(g=\psi_1(a)>0\), so the positive real square root used in the matrix formulas is
well defined.  Let
\[
 \mathcal V=\{(a,\lambda):a\in\R\setminus\{0,-1,-2,\ldots\},\ \lambda>0,\ Z(a,\lambda)\ne0\}.
\]
On each connected component of \(\mathcal V\), \(Z\) has constant sign, so
\(\partial_a\log|Z|=Z_a/Z\) and the corresponding mixed and second derivative
identities hold there.  Finally, \eqref{eq:I-Z} gives
\(\log|I_{a-1}(2\sqrt\lambda)|=\tfrac{a-1}{2}\log\lambda+\log|Z(a,\lambda)|\) where
\(Z\ne0\), leaving the derivations of \eqref{eq:ABC-log},
\eqref{eq:bessel-congruence}, and \eqref{eq:D-Delta} unchanged on that
component.
\end{proof}

Away from the poles the leading small-\(z\) term keeps its positive sign.  For
\(a=\nu+1\in\R\setminus\{0,-1,-2,\ldots\}\) the ascending series gives
\(I_\nu(z)=(z/2)^\nu\Gamma(\nu+1)^{-1}(1+O(z^2))\), so \(I_\nu\) has constant sign
near \(z=0\); by \Cref{lem:continuation}, \eqref{eq:Delta1-sharp} and
\eqref{eq:D-Delta} with \(Z(a,0)=1/\Gamma(a)\) and \(\lambda=z^2/4\) give
\begin{equation}\label{eq:small-z-D}
 D_\nu(z)=\frac{2\bigl(a\psi_1(a)-1\bigr)}{a^3\psi_1(a)}z^2+O(z^4)
 \qquad(z\downarrow0),
\end{equation}
whose coefficient is positive for every such \(a\): for \(a>0\) by
\eqref{eq:trig-lower}, and for negative nonintegral \(a\) because
\eqref{eq:trigamma-partial-fraction} makes \(a\psi_1(a)-1\) and \(a^3\psi_1(a)\) both
negative.  Coefficientwise positivity of \(\Delta\) nonetheless fails already on
the next interval where \(1/\Gamma(a)>0\): for \(-3<\nu<-2\) one has
\(a\in(-2,-1)\) and \(\Gamma(a)>0\).

\begin{proposition}[Degree-two obstruction beyond the positive-series domain]\label{prop:negative-coeff-failure}
If \(-2<a<-1\), equivalently \(-3<\nu<-2\), then
\begin{equation}\label{eq:negative-Delta2}
 [\lambda^2]\Delta(a,\lambda)=\Delta_2(a)<0.
\end{equation}
Thus the coefficientwise positivity theorem for \(\Delta\) cannot be extended to
that interval.
\end{proposition}

\begin{proof}
By \Cref{lem:continuation}, the identity \eqref{eq:Delta2-Q}--\eqref{eq:Qstar-decomp} persists on \(-2<a<-1\).  Set \(b=-a\in(1,2)\), so that \(a+1\in(-1,0)\).  The denominator in \eqref{eq:Delta2-Q} is then negative: \(a^6>0\) and \(\Gamma(a)^4>0\), while \((a+1)^3<0\).  It therefore suffices to show \(Q_*(a,t)>0\).  By \eqref{eq:trigamma-partial-fraction}, \(t=\psi_1(a+1)>0\) and
\[
 \mathcal R_a=t-\frac1{a+1}>\frac1{b-1}.
\]
The first term in \eqref{eq:Qstar-decomp} is positive.  The second and third terms satisfy
\begin{align*}
 &2b^2(b-1)^2(8b^2-3b+1)\mathcal R_a
 -2b(b-1)(5b-3)\\
 &\qquad>
 2b(b-1)\left[8b^3-3b^2-4b+3\right]>0.
\end{align*}
Indeed, the polynomial in brackets equals \(4\) at \(b=1\) and has derivative \(24b^2-6b-4>0\) for \(b\ge1\).  Hence \(Q_*(a,t)>0\), proving \eqref{eq:negative-Delta2}.
\end{proof}

A complete classification beyond \(\nu>-1\) appears to require separate
treatment.

\section{Comparison with related positivity results}\label{sec:context}

The scalar quantity \(A=Z_a^2-ZZ_{aa}\) is the infinitesimal symmetric-shift
Tur\'anian of the reciprocal-gamma series.  Indeed,
\(Z(a+\delta,\lambda)Z(a-\delta,\lambda)-Z(a,\lambda)^2=-A(a,\lambda)\delta^2+O(\delta^4)\).
The scalar coefficient sign of \(A\) is the infinitesimal case of the
coefficientwise reciprocal-gamma Tur\'an inequalities of Kalmykov and Karp
\cite[Thm.~3.1 and Ex.~1]{KalmykovKarp2013Reciprocal} (the integer-shift form
in \cite[Lem.~1 and Thm.~3]{KalmykovKarp2017}), whose gamma-ratio results
similarly treat scalar two-function Tur\'anians under parameter shifts
\cite{KalmykovKarp2013Ratios,KalmykovKarp2014,KarpSitnik2010}.  The coefficientwise thresholds of Karp and Zhang \cite[\S3]{KarpZhang2024} are set by the multiplicity of Pochhammer symbols in such scalar series; the threshold here is instead the deformation parameter \(\kappa\) of a coupled matrix determinant.

The additional quantities \(B\) and \(C\) are not scalar parameter-shift
Tur\'anians: \(B\) contains a mixed parameter--Euler derivative, while \(C\)
contains the sharp argument-direction expression and the trigamma
normalization that becomes the least argument-independent correction under the
Bessel congruence.  The result proved here is the coefficientwise positivity of the
coupled determinant \(AC-\psi_1(a)B^2\).

Classical Bessel Tur\'an inequalities compare neighboring orders, for example
\(I_\nu^2-I_{\nu-1}I_{\nu+1}\), or control ratios and logarithmic derivatives
\cite{Baricz2008,Baricz2010,BariczPonnusamy2013,Segura2011}; series
representations of the modified-Bessel Tur\'anian \cite[Thms.~1 and 5]{MezoBaricz2017},
positivity of Tur\'an determinants for symmetric orthogonal polynomials
\cite[Thm.~1 and Cor.~1]{BustozIsmail1997}, and higher-order Tur\'an and Hankel determinants in the
order \cite[\S\S2.2, 3]{BariczPogany2014} have also been studied, as has the absolute monotonicity in
the argument of Hankel determinants with special-function entries
\cite[Thm.~2.3]{IsmailLaforgia2007}.  In contrast, \eqref{eq:bessel-main} couples second
curvature in the order with first mixed order--argument curvature and an
argument-direction expression.  Derivatives of \(I_\nu\) with respect to the order have been
studied in their own right, with closed forms and integral representations for
\(\partial_\nu I_\nu\) \cite[\S10.38]{NISTDLMFOrder},\,\cite{GonzalezSantander2018}
and a closed form for \(\partial_\nu^2I_\nu\) at integer orders
\cite[eq.~(7.1)]{Brychkov2016}.  For \(J_\nu\) and \(Y_\nu\) the corresponding
integral representations, with uniform Airy asymptotics for large order, are
those of \cite[Props.~2.1 and 2.2, Thm.~4.2]{Dunster2017}; the
modified-function representations of \cite{GonzalezSantander2018} follow a
similar derivation.  The quantity \(P_\nu\) couples a first derivative in \(\nu\)
to the argument, whereas \(G_\nu\) is second order in \(\nu\) and is required
for every \(\nu>-1\), not only at integer orders.  Positivity of the leading entry \(G_\nu\) alone is the strict log-concavity
of \(\nu\mapsto I_\nu(z)\) on \((-1,\infty)\), which is known
\cite[Ex.~1]{KalmykovKarp2013Reciprocal}; a coefficient argument gives it on the
same range in the two-sided shifted form
\(1<I_\nu^2/(I_{\nu+\delta}I_{\nu-\delta})<\Gamma(\nu+\delta+1)\Gamma(\nu-\delta+1)/\Gamma(\nu+1)^2\)
for \(0<\delta<\nu+1\)
\cite[Thm.~3]{Nanthanasub2019}.  These results concern individual order
derivatives or scalar order log-concavity; neither the coupled determinant nor
the pointwise inequality of \Cref{thm:bessel}, with its argument-independent
correction \(4/\psi_1(\nu+1)\), appears to have an antecedent among them.

A recent preprint of Salazar establishes that the kernel \((z,\nu)\mapsto I_\nu(z)\) is strictly
totally positive for \(z>0\) and \(\nu\ge0\) \cite{Salazar2026Darboux}, extending
integer-order results of Buchstaber and Glutsyuk \cite{BuchstaberGlutsyuk2019};
a companion preprint builds all-order Hankel and Tur\'an
determinant hierarchies for special functions from positive moment
representations \cite{Salazar2026Hankel}.  Their determinant hierarchies are,
respectively, of evaluation-minor and Hankel-moment type; none of these papers
states or proves the coefficientwise mixed-derivative determinant studied here.
The signs of all total-positivity minors of \((z,\nu)\mapsto I_\nu(z)\), including
their confluent limits, are preserved under positive row and column rescalings
of the kernel, since a positive factor acts on a confluent jet matrix by a
triangular transformation with positive diagonal and multiplies each minor by a
positive number.  The Schur determinant of \eqref{eq:bessel-main} is not
preserved.  To make this explicit, write \(F(\nu,y)=\log I_\nu(e^y)\), so that
\(\partial_y=\Thetaz\) under \(z=e^y\) and \eqref{eq:Gnu}--\eqref{eq:Hnu} read
\[
 G_\nu=-F_{\nu\nu},\qquad
 P_\nu=F_{\nu y},\qquad
 H_\nu=2(F_y-\nu)-F_{yy}.
\]
Replacing \(I_\nu(e^y)\) by \(e^{cy}I_\nu(e^y)\) gives \(\widetilde F=F+cy\), with
\[
 \widetilde G_\nu=-\widetilde F_{\nu\nu}=G_\nu,
 \qquad
 \widetilde P_\nu=\widetilde F_{\nu y}=P_\nu,
 \qquad
 \widetilde H_\nu=2(\widetilde F_y-\nu)-\widetilde F_{yy}=H_\nu+2c.
\]
By \eqref{eq:Dnu-def} the rescaled kernel therefore has
\[
 \widetilde D_\nu(z)=D_\nu(z)+2cG_\nu(z).
\]
Here \(G_\nu=A/Z^2>0\), so at any fixed \(z_0>0\) every
\(c<-D_\nu(z_0)/\bigl(2G_\nu(z_0)\bigr)\) gives \(\widetilde D_\nu(z_0)<0\).  The
sign of \(D_\nu\) is therefore not invariant under positive rescaling of the
kernel, so \eqref{eq:bessel-main} is not a rescaling-covariant confluent
total-positivity minor.

\section{Further questions}\label{sec:questions}

\begin{question}\label{q:wright}
The series \(Z\) is the \(\rho=1\) member of a Wright-type family.  For \(\rho>0\), let
\[
 Z_\rho(a,\lambda)=\sum_{k\ge0}\frac{\lambda^k}{k!\,\Gamma(a+\rho k)},
\]
and, for \(\kappa\in\R\), define
\[
 \frac{A_\rho}{Z_\rho^2}=-\partial_a^2\log Z_\rho,
 \qquad
 \frac{B_\rho}{Z_\rho^2}=1+\rho\,\partial_a\Thetal\log Z_\rho,
\]
\[
 \frac{C_{\rho,\kappa}}{\psi_1(a)Z_\rho^2}
 =\frac1{\psi_1(a)}
 +\kappa\rho\,\Thetal\log Z_\rho
 -\rho^2\Thetal^2\log Z_\rho,
 \qquad
 \Delta_\rho^{(\kappa)}=A_\rho C_{\rho,\kappa}-\psi_1(a)B_\rho^2.
\]
For which pairs \((\rho,\kappa)\) does \(\Delta_\rho^{(\kappa)}\) have strictly positive Maclaurin coefficients in every positive degree for all \(a>0\); and for each \(\rho>0\), what is the sharp threshold in \(\kappa\)?  At \(\rho=1\) the present theorem gives exactly \(\kappa\ge1\), and its proof makes essential use of the exact \(\rho=1\) Chu--Vandermonde convolution.
\end{question}

\begin{question}\label{q:tp-origin}
Is there a canonical normalization, bordered kernel, or nonlocal minor
associated with \((z,\nu)\mapsto I_\nu(z)\) whose positivity is equivalent to
the Bessel--Schur inequality of \Cref{thm:bessel}?

The rescaling computation of \Cref{sec:context} shows that any such derivation
must incorporate additional Bessel-specific normalization or nonlocal
information.  A related question is whether the coefficientwise positivity of
\(AC-\psi_1(a)B^2\) follows from an appropriately strengthened form of scalar
reciprocal-gamma log-concavity.
\end{question}

\section*{Declarations}

\noindent\textbf{Funding.} This research received no external funding.

\smallskip
\noindent\textbf{Competing interests.} The author has no relevant interests to disclose.

\smallskip
\noindent\textbf{Data availability.} No datasets were generated or analysed.

\smallskip
\noindent\textbf{Computational verification.} The proofs do not depend on numerical computation.  Scripts used for verification and figures are available at \url{https://github.com/johnnyshields/papers}.

\smallskip
\noindent\textbf{Generative AI use.} During the preparation of this work, the author used OpenAI's ChatGPT (GPT-5.6 Sol) and Anthropic's Claude (Opus 5, Opus 4.8) to assist with manuscript drafting and the development of auxiliary verification code.  After using these tools, the author reviewed and edited as needed.  The author takes full responsibility for the content of the work.

\begin{thebibliography}{99}

\bibitem{Baricz2008}
\'{A}. Baricz,
\emph{Functional inequalities involving Bessel and modified Bessel functions of the first kind},
Expo. Math. \textbf{26} (2008), 279--293,\newline DOI: \href{https://doi.org/10.1016/j.exmath.2008.01.001}{10.1016/j.exmath.2008.01.001}.

\bibitem{Baricz2010}
\'{A}. Baricz,
\emph{Tur\'{a}n type inequalities for modified Bessel functions},
Bull. Aust. Math. Soc. \textbf{82} (2010), 254--264,\newline DOI: \href{https://doi.org/10.1017/S000497271000002X}{10.1017/S000497271000002X}.

\bibitem{BariczPogany2014}
\'{A}. Baricz and T. K. Pog\'{a}ny,
\emph{Tur\'{a}n determinants of Bessel functions},
Forum Math. \textbf{26} (2014), 295--322,\newline DOI: \href{https://doi.org/10.1515/form.2011.160}{10.1515/form.2011.160}.

\bibitem{BariczPonnusamy2013}
\'{A}. Baricz and S. Ponnusamy,
\emph{On Tur\'{a}n type inequalities for modified Bessel functions},
Proc. Amer. Math. Soc. \textbf{141} (2013), 523--532,\newline DOI: \href{https://doi.org/10.1090/S0002-9939-2012-11325-5}{10.1090/S0002-9939-2012-11325-5}.

\bibitem{Brychkov2016}
Yu.~A. Brychkov,
\emph{Higher derivatives of the Bessel functions with respect to the order},
Integral Transforms Spec. Funct. \textbf{27} (2016), 566--577,\newline DOI: \href{https://doi.org/10.1080/10652469.2016.1164156}{10.1080/10652469.2016.1164156}.

\bibitem{BuchstaberGlutsyuk2019}
V. M. Buchstaber and A. A. Glutsyuk,
\emph{Total positivity, Grassmannian and modified Bessel functions},
in \emph{Functional Analysis and Geometry: Selim Grigorievich Krein Centennial},
Contemp. Math. \textbf{733}, Amer. Math. Soc., 2019, 97--107,\newline DOI: \href{https://doi.org/10.1090/conm/733/14736}{10.1090/conm/733/14736}.

\bibitem{BustozIsmail1997}
J.~Bustoz and M.~E.~H. Ismail,
\emph{Tur\'{a}n inequalities for symmetric orthogonal polynomials},
Int. J. Math. Math. Sci. \textbf{20} (1997), 1--7,\newline DOI: \href{https://doi.org/10.1155/S016117129700001X}{10.1155/S016117129700001X}.

\bibitem{Dunster2017}
T.~M. Dunster,
\emph{On the order derivatives of Bessel functions},
Constr. Approx. \textbf{46} (2017), 47--68,\newline DOI: \href{https://doi.org/10.1007/s00365-016-9355-1}{10.1007/s00365-016-9355-1}.

\bibitem{GonzalezSantander2018}
J.~L. Gonz\'alez-Santander,
\emph{Closed-form expressions for derivatives of Bessel functions with respect to the order},
J. Math. Anal. Appl. \textbf{466} (2018), 1060--1081,\newline DOI: \href{https://doi.org/10.1016/j.jmaa.2018.06.043}{10.1016/j.jmaa.2018.06.043}.

\bibitem{IsmailLaforgia2007}
M.~E.~H. Ismail and A.~Laforgia,
\emph{Monotonicity properties of determinants of special functions},
Constr. Approx. \textbf{26} (2007), 1--9,\newline DOI: \href{https://doi.org/10.1007/s00365-005-0627-4}{10.1007/s00365-005-0627-4}.

\bibitem{KalmykovKarp2013Reciprocal}
S. I. Kalmykov and D. B. Karp,
\emph{Log-concavity for series in reciprocal gamma functions and applications},
Integral Transforms Spec. Funct. \textbf{24} (2013), 859--872,\newline DOI: \href{https://doi.org/10.1080/10652469.2013.764874}{10.1080/10652469.2013.764874}.

\bibitem{KalmykovKarp2013Ratios}
S. I. Kalmykov and D. B. Karp,
\emph{Log-convexity and log-concavity for series in gamma ratios and applications},
J. Math. Anal. Appl. \textbf{406} (2013), 400--418,\newline DOI: \href{https://doi.org/10.1016/j.jmaa.2013.04.061}{10.1016/j.jmaa.2013.04.061}.

\bibitem{KalmykovKarp2014}
S. I. Kalmykov and D. B. Karp,
\emph{Logarithmic concavity of series in gamma ratios},
Russian Math. \textbf{58} (2014), 63--68,\newline DOI: \href{https://doi.org/10.3103/S1066369X14060073}{10.3103/S1066369X14060073}.

\bibitem{KalmykovKarp2017}
S. I. Kalmykov and D. B. Karp,
\emph{Log-concavity and Tur\'{a}n-type inequalities for the generalized hypergeometric function},
Anal. Math. \textbf{43} (2017), 567--580,\newline DOI: \href{https://doi.org/10.1007/s10476-017-0503-z}{10.1007/s10476-017-0503-z}.

\bibitem{KarpSitnik2010}
D.~Karp and S.~M. Sitnik,
\emph{Log-convexity and log-concavity of hypergeometric-like functions},
J. Math. Anal. Appl. \textbf{364} (2010), 384--394,\newline DOI: \href{https://doi.org/10.1016/j.jmaa.2009.10.057}{10.1016/j.jmaa.2009.10.057}.

\bibitem{KarpZhang2024}
D.~Karp and Y.~Zhang,
\emph{Log-concavity and log-convexity of series containing multiple Pochhammer symbols},
Fract. Calc. Appl. Anal. \textbf{27} (2024), 458--486,\newline DOI: \href{https://doi.org/10.1007/s13540-023-00238-0}{10.1007/s13540-023-00238-0}.

\bibitem{MezoBaricz2017}
I.~Mez\H{o} and \'{A}.~Baricz,
\emph{Properties of the Tur\'{a}nian of modified Bessel functions},
Math. Inequal. Appl. \textbf{20} (2017), 991--1001,\newline DOI: \href{https://doi.org/10.7153/mia-2017-20-63}{10.7153/mia-2017-20-63}.

\bibitem{Nanthanasub2019}
T.~Nanthanasub, B.~Novaprateep and N.~Wichailukkana,
\emph{The logarithmic concavity of modified Bessel functions of the first kind and its related functions},
Adv. Difference Equ. \textbf{2019} (2019), art.~379, 14 pp.,\newline DOI: \href{https://doi.org/10.1186/s13662-019-2309-8}{10.1186/s13662-019-2309-8}.

\bibitem{NISTDLMFBessel}
NIST Digital Library of Mathematical Functions,
\emph{Chapter 10: Bessel Functions},
\url{https://dlmf.nist.gov/10}, accessed 22 July 2026.

\bibitem{NISTDLMFGamma}
NIST Digital Library of Mathematical Functions,
\emph{Chapter 5: Gamma Function},
\url{https://dlmf.nist.gov/5}, accessed 19 July 2026.

\bibitem{NISTDLMFOrder}
NIST Digital Library of Mathematical Functions,
\emph{Section 10.38: Derivatives with Respect to Order},
\url{https://dlmf.nist.gov/10.38}, accessed 22 July 2026.

\bibitem{Salazar2026Darboux}
D. S. P. Salazar,
\emph{Strict total positivity from spectral Darboux and Toeplitz smoothing mechanisms},
arXiv:2607.02778, 2026,\newline DOI: \href{https://doi.org/10.48550/arXiv.2607.02778}{10.48550/arXiv.2607.02778}.

\bibitem{Salazar2026Hankel}
D. S. P. Salazar,
\emph{Order-moment transport and Hankel determinants in special-function inequalities},
arXiv:2606.31647, 2026,\newline DOI: \href{https://doi.org/10.48550/arXiv.2606.31647}{10.48550/arXiv.2606.31647}.

\bibitem{Segura2011}
J. Segura,
\emph{Bounds for ratios of modified Bessel functions and associated Tur\'{a}n-type inequalities},
J. Math. Anal. Appl. \textbf{374} (2011), 516--528,\newline DOI: \href{https://doi.org/10.1016/j.jmaa.2010.09.030}{10.1016/j.jmaa.2010.09.030}.

\bibitem{Watson1944}
G. N. Watson,
\emph{A Treatise on the Theory of Bessel Functions}, 2nd ed.,
Cambridge University Press, 1944.

\end{thebibliography}

\end{document}
